% META: {"id": "1NSI-WEB-ME-002", "chapitre": "1NSI-WEB-IHM", "type_objet": "methode", "capacites": ["P-IHM-02"], "capacites_codes": ["C3"], "methodes": ["M2"], "mode_creation": "ex_nihilo", "sources_inspiration": [], "status": "needs_review"}
% BEGIN-VERIFY
% gestionnaires = {}
% def ajouter_gestionnaire(id_composant, fonction):
%     gestionnaires[id_composant] = fonction
% def declencher_clic(id_composant):
%     if id_composant in gestionnaires:
%         return gestionnaires[id_composant]()
%     return None
% panier = {"articles": 0}
% def ajouter_article():
%     panier["articles"] = panier["articles"] + 1
%     return panier["articles"]
% ajouter_gestionnaire("bouton_ajouter", ajouter_article)
% assert declencher_clic("bouton_ajouter") == 1
% assert declencher_clic("bouton_ajouter") == 2
% assert panier["articles"] == 2
% assert declencher_clic("bouton_inconnu") is None
% END-VERIFY
\begin{fichemethode}{M2}{Lire et modifier la fonction exécutée lors d'un clic}

\textbf{Quand l'utiliser ?} Quand l'énoncé demande ce qu'affiche une page après un ou
plusieurs clics, ou de \textbf{modifier} le comportement d'un bouton.

\textbf{Pas à pas :}
\begin{enumerate}
  \item \textbf{Retrouve l'association} entre le composant et la fonction : cherche l'appel à
        \lstinline{ajouter_gestionnaire}. Il relie un \emph{identifiant} à une \emph{fonction}.
        Tant qu'aucune association n'existe, cliquer ne produit rien.
  \item \textbf{Lis la fonction, pas son nom.} Un nom rassurant ne prouve rien : c'est le corps
        qui décide de l'effet.
  \item \textbf{Repère ce qui est modifié durablement.} Une variable modifiée dans la fonction
        (souvent un dictionnaire d'état) \textbf{conserve} sa valeur d'un clic au suivant.
        C'est ce qui rend les clics successifs différents les uns des autres.
  \item \textbf{Simule les clics un par un}, en écrivant l'état après chacun. Ne saute pas
        d'étape : c'est là que les erreurs apparaissent.
  \item \textbf{Pour modifier le comportement}, change le corps de la fonction ou associe une
        autre fonction au même identifiant. Ne touche pas au HTML : le bouton, lui, ne change pas.
\end{enumerate}

\textbf{Exemple rédigé.} Un bouton doit incrémenter le nombre d'articles du panier.

\begin{python}
panier = {"articles": 0}

def ajouter_article():
    panier["articles"] = panier["articles"] + 1
    return panier["articles"]

ajouter_gestionnaire("bouton_ajouter", ajouter_article)
\end{python}

Premier clic : \lstinline{declencher_clic("bouton_ajouter")} renvoie \lstinline{1}.
Deuxième clic : il renvoie \lstinline{2}, car \lstinline{panier} a gardé sa valeur.
Un clic sur un identifiant sans gestionnaire renvoie \lstinline{None} : rien ne se produit.

\textbf{Erreur fréquente.} Écrire \lstinline{ajouter_gestionnaire("bouton", ajouter_article())}
avec des parenthèses. On appelle alors la fonction \emph{immédiatement} et on associe son
\emph{résultat} au bouton. Il faut transmettre la fonction elle-même, sans parenthèses.

\end{fichemethode}
